\documentclass[12pt,a4paper]{report}

% ==========================================================
%                    PACKAGES
% ==========================================================

\usepackage[utf8]{inputenc}     % Encodage
\usepackage[T1]{fontenc}        % Accents
\usepackage[french]{babel}      % Langue française

\usepackage{amsmath}            % Math
\usepackage{amsfonts}
\usepackage{amssymb}

\usepackage{graphicx}           % Figures
\usepackage{hyperref}           % Liens cliquables
\usepackage{geometry}           % Marges
\usepackage{siunitx}            % Unités
\usepackage{caption}
\usepackage{subcaption}

\geometry{margin=2.5cm}

% ==========================================================
%                    DOCUMENT
% ==========================================================

\begin{document}

\title{Étude des performances électrochimiques des cathodes lithium-ion par DFT}
\author{Hajar Motahhir}
\date{\today}

\maketitle

\tableofcontents

% ==========================================================
%                CHAPITRE 1 : INTRODUCTION
% ==========================================================

\chapter{Introduction générale}
\label{ch:introduction}


Les besoins énergétiques mondiaux augmentent fortement en raison de l’industrialisation,
de la transition énergétique et du développement des transports électriques.
Les systèmes de stockage d’énergie deviennent ainsi essentiels pour assurer
une exploitation efficace des ressources énergétiques.

Les batteries lithium-ion (LIBs) sont actuellement les dispositifs de stockage les plus utilisés,
en raison de leur haute densité énergétique, leur efficacité et leur longue durée de vie.

La performance globale des LIBs dépend principalement de la nature des matériaux cathodiques,
qui contrôlent la capacité, la tension, la stabilité structurale et les mécanismes électrochimiques.
Cependant, les matériaux cathodiques classiques présentent encore des limitations :
\begin{itemize}
    \item faible stabilité à haute tension,
    \item dégradation structurale,
    \item perte de capacité,
    \item coût élevé,
    \item ou encore mécanismes redox complexes.
\end{itemize}

La conception de nouvelles cathodes à haute performance nécessite donc une compréhension profonde
des propriétés structurales et électroniques à l’échelle atomique.
Les approches théoriques, et en particulier la \emph{Density Functional Theory} (DFT),
jouent un rôle essentiel pour explorer des matériaux,
comparer différentes compositions et prédire leurs performances avant leur synthèse expérimentale.

\section{Motivation}

Les cathodes comme LiCoO$_2$, LiNiMnCoO$_2$ ou LiFePO$_4$ 
sont largement utilisées mais souffrent encore de limitations.
Le dopage ou la substitution dans les structures oxydes permet :
\begin{itemize}
    \item d’améliorer la stabilité,
    \item d’augmenter la conductivité,
    \item de modifier les mécanismes redox,
    \item ou encore de réduire les barrières de diffusion.
\end{itemize}

La DFT fournit des outils robustes pour analyser ces effets,
notamment les états électroniques, les mécanismes redox,
les énergies de formation ou les tensions électrochimiques.

\section{Problématique}

Malgré les progrès réalisés,
les mécanismes électroniques dans les cathodes,
notamment les contributions métal–oxygène,
restent encore mal compris.
La présence de corrélations électroniques fortes dans les métaux de transition
rend les analyses expérimentales complexes.

La problématique générale de ce travail est donc :

\begin{quote}
\textit{Comment utiliser les calculs DFT pour prédire les performances électrochimiques 
des cathodes lithium-ion et comprendre les effets du dopage sur leurs propriétés structurales et électroniques ?}
\end{quote}

\section{Objectifs}

Les objectifs principaux de cette étude sont :
\begin{enumerate}
    \item Étudier les structures cristallines des cathodes,
    \item Calculer les propriétés électroniques (bandes, DOS, PDOS),
    \item Analyser les mécanismes redox,
    \item Estimer les barrières de diffusion Li$^+$,
    \item Calculer les tensions théoriques,
    \item Étudier les effets du dopage (C, N, Si, P),
    \item Identifier les matériaux prometteurs pour les LIBs.
\end{enumerate}

\section{Méthodologie générale}

La méthodologie repose sur :
\begin{itemize}
    \item construction des modèles cristallins,
    \item optimisation structurale,
    \item calcul des états électroniques,
    \item simulation de la dé-lithiation,
    \item calcul des énergies,
    \item analyse des mécanismes électrochimiques,
    \item comparaison entre matériaux dopés et non dopés.
\end{itemize}

Les calculs sont réalisés à l’aide de codes ab initio comme WIEN2k ou VASP,
utilisant des fonctionnelles GGA, GGA+U ou hybrides selon les besoins.

\section{Organisation de la thèse}

La thèse est structurée comme suit :

\begin{description}
    \item[Chapitre 1 :] Introduction générale.
    \item[Chapitre 2 :] Bases théoriques de la DFT et méthodologie.
    \item[Chapitre 3 :] Étude des cathodes non dopés.
    \item[Chapitre 4 :] Effet du dopage sur les propriétés structurales et électroniques.
    \item[Chapitre 5 :] Analyse électrochimique et mécanismes redox.
    \item[Chapitre 6 :] Conclusions et perspectives.
\end{description}

\end{document}
